\chapter{The Language of Mathematics}\label{ch:language}

\begin{center}\itshape
The equation $x^2=1$ has a solution. It has a unique solution. Every real number solves it.\\
Three sentences, the same equation, three different mathematical questions.
\end{center}

Before solving an equation, we need to know what its question means. Over the real numbers,
$x^2=1$ has the two solutions $1$ and $-1$. Over the positive real numbers it has just one.
The calculation has hardly changed; the statement has. The domain, the words ``there is''
and ``every'', and the distinction between existence and uniqueness carry mathematical information.
Learning to read that information is the first step towards writing proofs.

Here is a second question. Someone claims that every increasing function is positive, and
offers the examples $x\mapsto x+1$ and $x\mapsto x+10$ for $x\ge0$. How much evidence is enough?
None of these examples addresses the function $x\mapsto x-1$ on the same domain. A single
counterexample defeats the claim. A thousand favourable examples would not prove it.
This asymmetry will recur throughout the book: experimentation helps us find a statement;
a proof must explain why all the cases covered by its words obey it.

We begin with familiar arithmetic and build a working language for unfamiliar mathematics.
The point of the symbols is to make our obligations visible: what is given, what may be
chosen, what must be established, and which choices may depend on which others.

\section{Definitions, assumptions and conclusions}

A \emph{definition} gives a word an exact meaning. An integer $n$ is \emph{even} if
$n=2k$ for some integer $k$. This definition is a test, not a theorem awaiting proof.
To prove that $18$ is even, produce the integer $k=9$. To use the assumption that
$n$ is even, write $n=2k$ for some integer $k$ and keep track of that integer.
Different words can describe the same objects, but that agreement then needs a proof.

An \emph{axiom} is an assumption that specifies the mathematical setting. A
\emph{theorem} is a statement proved from the assumptions and earlier results.
The words \emph{proposition}, \emph{lemma} and \emph{corollary} also denote proved statements:
a lemma serves another proof; a corollary follows readily from an earlier result.
These names describe the statement's role, not different standards of correctness.
A \emph{conjecture} is a proposed statement for which a proof or disproof is still sought.
A conjecture is not available as a theorem merely because many examples support it.

\begin{reminder}[our starting point]
We use ordinary classical logic and the usual arithmetic and order of natural numbers
$\N=\{0,1,2,\dots\}$, integers $\Z$, rational numbers $\Q$, and real numbers $\R$.
A rational number has the form $a/b$ with $a,b\in\Z$ and $b\ne0$.
The real-number axioms include completeness, whose consequences are developed in the
sequences chapter. Well ordering of $\N$ is stated and used in the next chapter.
We use sets and functions in their ordinary mathematical sense, rather than construct
set theory from formal axioms. Later chapters state explicitly the elementary-function
and integration facts they assume. Nothing here requires a previous course in proofs.
\end{reminder}

In a theorem of the form ``if $H$, then $C$'', $H$ is the \emph{hypothesis} and $C$ the
\emph{conclusion}. A proof may use $H$, the definitions, the stated foundations and results
already established. It may not use $C$ as an extra assumption and then announce that $C$
has been proved. Working backwards from $C$ can be useful on scrap paper, but the final
argument must justify every step from the available information.

\begin{examplex}[unpacking a definition]
The square of every even integer is even. Indeed, let $n$ be an even integer. By definition
$n=2k$ for an integer $k$, so $n^2=4k^2=2(2k^2)$. Since $2k^2$ is an integer, the definition
says that $n^2$ is even. The witness for the hypothesis was $k$; the witness for the
conclusion is $2k^2$. The proof is a calculation guided by two uses of one definition.
\end{examplex}

\section{Statements, connectives and implication}

A \emph{statement} is a sentence with a definite truth value in its stated setting.
``$7>3$'' is true; ``$7<3$'' is false. The expression ``$x>3$'' becomes a statement
after $x$ is specified or quantified. We call it a \emph{predicate} depending on $x$.
An equation such as $1/x=2$ is only meaningful where $x\ne0$. Establish the domain
before treating an expression as a predicate.

For statements $P$ and $Q$, the connective $P\land Q$ means ``$P$ and $Q$'', and is
true exactly when both are true. The connective $P\lor Q$ means ``$P$ or $Q$'', and
is true when at least one is true, including when both are true. This is the inclusive
mathematical meaning of ``or''. The negation $\neg P$ means ``$P$ is false''.
In prose we usually write these words rather than the symbols.

The implication $P\Rightarrow Q$ says that $P$ cannot be true while $Q$ is false.
Equivalently, it has the same truth value as $(\neg P)\lor Q$.
Its four possible cases are worth seeing once:
\[
\begin{array}{c|c|c}
P&Q&P\Rightarrow Q\\\hline
\text{true}&\text{true}&\text{true}\\
\text{true}&\text{false}&\text{false}\\
\text{false}&\text{true}&\text{true}\\
\text{false}&\text{false}&\text{true}
\end{array}
\]
An implication promises something when its hypothesis holds; it promises nothing
about cases where that hypothesis fails. Thus ``if $x>2$, then $x^2>4$'' is true
for every real $x$, including $x=0$. At $x=0$ the hypothesis is false, so it does
not provide a counterexample. This convention is called \emph{vacuous truth}.

An implication does not assert a causal connection. It also does not permit us to
reverse the arrow. From $x=1\Rightarrow x^2=1$, the converse $x^2=1\Rightarrow x=1$
does not follow: $x=-1$ defeats it. Nor does a false hypothesis force a false conclusion.
From ``if $n$ is a multiple of $4$, then $n$ is even'', we cannot deduce that every integer
which is not a multiple of $4$ is not even; $6$ is an immediate test.
Here ``$a$ is a multiple of $b$'' means $a=bk$ for an integer $k$.

The equivalence $P\Leftrightarrow Q$ means that both $P\Rightarrow Q$ and $Q\Rightarrow P$
hold. It is read ``$P$ if and only if $Q$'', often abbreviated to ``iff''.
For real $x$, $x^2=0\Leftrightarrow x=0$: one direction substitutes $0$; in the other,
if $x\ne0$, division by $x$ would give $x=0$, a contradiction.

If $P\Rightarrow Q$, then $P$ is a \emph{sufficient} condition for $Q$, and $Q$ is a
\emph{necessary} condition for $P$. Being a multiple of $4$ is sufficient for being even;
being even is necessary for being a multiple of $4$. Neither phrase reverses the implication.
To claim a necessary and sufficient condition is to claim an equivalence, so two directions
must be justified.

\begin{think}
For a real number $x$, is $x>1$ necessary for $x^2>1$? Is it sufficient?
Test a negative number before deciding. Which necessary and sufficient condition
for $x^2>1$ uses both positive and negative numbers?
\end{think}

\section{Quantifiers: who chooses what?}

The \emph{universal quantifier} $\forall$ means ``for every''; the \emph{existential
quantifier} $\exists$ means ``there exists''. The domain belongs to the statement:
\[
\forall x\in\R,\quad x^2\ge0;
\qquad
\exists x\in\R,\quad x^2=2.
\]
The first requires a proof covering an arbitrary real $x$. The second requires at
least one witness. For instance the usual positive square root of $2$ is a real witness;
its existence is a real-number fact developed from completeness later in the book.
Requiring a rational witness is a different question: a real square root does not
settle it. For an immediate example where the domain changes the answer,
$\exists x\in\R,\ 2x=1$ is true with $x=1/2$, whereas
$\exists x\in\Z,\ 2x=1$ is false, since twice an integer cannot lie strictly between $0$ and $2$.

To see how witnesses work without any existence theorem, compare
\[
\forall x\in\R\ \exists y\in\R:\ x+y=0
\quad\hbox{and}\quad
\exists y\in\R\ \forall x\in\R:\ x+y=0.
\]
The first is true: after $x$ is given, choose $y=-x$. This $y$ is allowed to depend on $x$.
The second is false: a single $y$ would have to satisfy both $0+y=0$ and $1+y=0$.
Those equations require $y=0$ and $y=-1$ simultaneously. The order of the quantifiers
specifies which information is available when a choice is made.

Adjacent universal quantifiers can be exchanged: ``for every $x$, for every $y$''
requires all pairs in either order. Adjacent existential quantifiers can also be
exchanged: a successful pair supplies both witnesses. Mixed quantifiers generally
cannot. When reading a long statement, mark the variables chosen universally and
draw an arrow from each existential choice to the earlier variables on which it may depend.
Later, the order $\forall\varepsilon>0\ \exists N\ \forall n\ge N$ will make the
definition of convergence precise: $N$ may depend on $\varepsilon$, but not on the
later index $n$.

The \emph{scope} of a quantifier is the part of the sentence to which it applies.
Use punctuation or parentheses when scope could be ambiguous. The statement
$\forall x\in\R,\ (x>0\Rightarrow x^2>0)$ quantifies an implication; it does not
say that every real number is positive. Quantifiers over an empty domain obey the
same conventions as implication: ``every element of the empty set has property $P$''
is true, whereas ``some element of the empty set has property $P$'' is false.
There is no element witnessing a failure of the first or a success of the second.

\subsection*{Negation without changing the question}

Negation reverses a quantifier and negates its predicate:
\[
\begin{aligned}
\neg(\forall x\in E,\ P(x))&\ \Longleftrightarrow\
\exists x\in E,\ \neg P(x),\\
\neg(\exists x\in E,\ P(x))&\ \Longleftrightarrow\
\forall x\in E,\ \neg P(x).
\end{aligned}
\]
The first says that a universal claim fails precisely when it has a counterexample.
The second says that no witness exists precisely when every candidate fails.
The familiar connective rules follow from the same meanings:
\[
\begin{aligned}
\neg(P\land Q)&\ \Longleftrightarrow\ (\neg P)\lor(\neg Q),\\
\neg(P\lor Q)&\ \Longleftrightarrow\ (\neg P)\land(\neg Q),\\
\neg(P\Rightarrow Q)&\ \Longleftrightarrow\ P\land(\neg Q).
\end{aligned}
\]
For example, ``$x$ does not satisfy $0<x<1$'' means $x\le0$ or $x\ge1$.
It does not mean $x<0$ and $x>1$. Endpoints belong to the negation.

Negating from the outside in gives
\[
\neg\bigl(\forall x\in\R\ \exists y\in\R:\ y>x\bigr)
\ \Longleftrightarrow\
\exists x\in\R\ \forall y\in\R:\ y\le x.
\]
The original statement is true, with witness $y=x+1$; its negation asserts the
existence of a greatest real number. Notice that negation preserves quantifier order
while changing each quantifier's type. It does not simply insert a minus sign into a formula.

\section{Sets: making the domain visible}

A \emph{set} is a collection of distinct objects, its \emph{elements}. We write $x\in A$
for membership and $x\notin A$ for non-membership. Listing braces ignore order and
repetition: $\{1,2,2\}=\{2,1\}$. The empty set $\varnothing$ has no elements.
Set-builder notation records a condition, as in
$\{x\in\R:x^2=1\}=\{-1,1\}$. The colon means ``such that''.

We write $A\subseteq B$ when every element of $A$ belongs to $B$:
\[
A\subseteq B\quad\Longleftrightarrow\quad
\forall x,\ (x\in A\Rightarrow x\in B).
\]
Equality of sets means equality of membership. Thus $A=B$ is equivalent to the two
inclusions $A\subseteq B$ and $B\subseteq A$. A \emph{proper} subset $A\subsetneq B$
is a subset unequal to $B$. Some authors use $\subset$ for either meaning; here we
write $\subseteq$ when equality is allowed. For every set $A$, $\varnothing\subseteq A$:
the implication defining this inclusion has no counterexample.

The \emph{union}, \emph{intersection} and \emph{difference} are
\[
\begin{aligned}
A\cup B&=\{x:x\in A\ \hbox{or}\ x\in B\},\\
A\cap B&=\{x:x\in A\ \hbox{and}\ x\in B\},\\
A\setminus B&=\{x:x\in A\ \hbox{and}\ x\notin B\}.
\end{aligned}
\]
These descriptions take their elements from the sets already under discussion.
For a complement we fix an \emph{ambient set} $U$ containing $A$ and write
$A^c=U\setminus A$. The complement of $\{1,2\}$ in $\{1,2,3\}$ is $\{3\}$;
its complement in $\N$ is a different set. We do not assume a set of all mathematical objects.

For a family of sets $(A_i)_{i\in I}$, membership in $\bigcup_{i\in I}A_i$ means
membership in at least one $A_i$; membership in $\bigcap_{i\in I}A_i$ means membership
in every $A_i$. When the intersection is taken inside a fixed $U$, the empty union
is $\varnothing$ and the empty intersection is $U$. These conventions match
existential and universal quantification over the empty index set.

\begin{proposition}[set operations and logic]\label{prop:lang-sets}
For subsets $A,B,C$ of a fixed set $U$,
\[
\begin{aligned}
A\cap(B\cup C)&=(A\cap B)\cup(A\cap C),\\
(A\cup B)^c&=A^c\cap B^c,\qquad
(A\cap B)^c=A^c\cup B^c.
\end{aligned}
\]
\end{proposition}
\begin{proof}
An element $x$ belongs to $A\cap(B\cup C)$ exactly when $x\in A$ and at least
one of $x\in B$, $x\in C$ holds. This is exactly membership in $(A\cap B)\cup(A\cap C)$:
if $x$ is in the first set, use whichever of $B,C$ contains it; in the reverse direction,
membership in either intersection supplies $x\in A$ and membership in the union.
For the second identity, an $x\in U$ is outside $A\cup B$ exactly when it is outside
both $A$ and $B$. For the third, $x\in U$ fails to belong to both $A$ and $B$
exactly when at least one membership fails. These are the connective negation rules.
\end{proof}

\subsection*{Power sets and ordered pairs}

The \emph{power set} $\mathcal P(A)$ is the set of all subsets of $A$. For example,
\[
\mathcal P(\{a,b\})=\{\varnothing,\{a\},\{b\},\{a,b\}\}.
\]
Membership and inclusion now coexist: $\{a\}\in\mathcal P(A)$ means the same as
$\{a\}\subseteq A$. The set $\varnothing$ differs from $\{\varnothing\}$:
the first has no elements; the second has one, namely the empty set.

An \emph{ordered pair} $(a,b)$ records a first and a second coordinate, so that
$(a,b)=(c,d)$ exactly when $a=c$ and $b=d$. Unlike a two-element set, an ordered
pair remembers order and allows repeated coordinates. The \emph{Cartesian product}
is $A\times B=\{(a,b):a\in A,\ b\in B\}$.
Thus $\{0,1\}\times\{2\}=\{(0,2),(1,2)\}$, whereas
$\{2\}\times\{0,1\}=\{(2,0),(2,1)\}$.
Repeated products give tuples $(x_1,\dots,x_n)$; in particular $\R^2$ is the coordinate
plane and $\R^n$ the set of real $n$-tuples. Products will describe successive
choices in counting and successive trials in probability.

\section{Functions: the rule and its two sets}

A \emph{function} $f:E\to F$ assigns to every element $x$ of its \emph{domain} $E$
exactly one element $f(x)$ of its \emph{codomain} $F$. The symbol $x\mapsto f(x)$
describes the assignment; $f$ is the function and $f(x)$ one value.
One can identify $f$ with its \emph{graph}, the subset
$\{(x,f(x)):x\in E\}$ of $E\times F$, together with its specified domain and codomain.
The graph meets each vertical fibre $\{x\}\times F$ in exactly one point.

A formula alone need not specify a function. The formula $1/x$ cannot define
$\R\to\R$ because it has no value at $0$; it does define $\R\setminus\{0\}\to\R$.
The notation ``$f(x)=\pm\sqrt x$'' does not select one value for positive $x$.
The formula $x^2$ does define a function on all of $\R$, and also on the smaller
domain $[0,\infty)=\{x\in\R:x\ge0\}$. Here and below $[a,b]$ includes its
endpoints, $(a,b)$ excludes them, and a square bracket at just one end includes
only that endpoint. Infinite endpoints are always excluded.

For $A\subseteq E$, the \emph{image} of $A$ is
\[
f(A)=\{f(x):x\in A\}
=\{y\in F:\exists x\in A,\ f(x)=y\}.
\]
The image of the whole domain, $f(E)$, is also called the \emph{range}. It may be
smaller than the codomain. For $B\subseteq F$, the \emph{inverse image}, or
\emph{preimage}, is
\[
f^{-1}(B)=\{x\in E:f(x)\in B\}.
\]
This notation makes sense for every function. It does not claim that $f$ has an
inverse function. For $f:\R\to\R$, $f(x)=x^2$, we have
$f(\{-2,1\})=\{4,1\}$, $f^{-1}(\{1\})=\{-1,1\}$ and
$f^{-1}((-\infty,0))=\varnothing$.

\begin{proposition}[inverse images preserve set operations]\label{prop:lang-preimages}
If $f:E\to F$ and $B,C\subseteq F$, then
\[
\begin{aligned}
f^{-1}(B\cup C)&=f^{-1}(B)\cup f^{-1}(C),\\
f^{-1}(B\cap C)&=f^{-1}(B)\cap f^{-1}(C),\\
f^{-1}(F\setminus B)&=E\setminus f^{-1}(B).
\end{aligned}
\]
\end{proposition}
\begin{proof}
For $x\in E$, membership in the left side of the first equality means that
$f(x)\in B$ or $f(x)\in C$, which is membership in the right side. Replace ``or''
by ``and'' for the second. For the third, $f(x)\notin B$ says exactly that $x$
is outside $f^{-1}(B)$, with complement taken in $E$.
\end{proof}

These equalities explain why events described by a random variable will be inverse
images. ``The measured value lies in $B$'' is a condition on the original outcomes.
An image instead collects values and can forget which inputs produced them.
That distinction matters: images do not always preserve intersections, as the workshop will show.

\begin{figure}[ht]
\centering
\fig{ch00-fibres}
\caption{A function can identify two inputs and miss a codomain value. Here
$f(x_1)=f(x_2)=a$, $f(x_3)=b$, and $c$ has no preimage. The dashed set is
one fibre; the image $\{a,b\}$ differs from the codomain $F=\{a,b,c\}$.}
\label{fig:language-fibres}
\end{figure}

\section{Injectivity, surjectivity and inverse functions}

A function $f:E\to F$ is \emph{injective} if equal outputs force equal inputs:
\[
\forall x,x'\in E,\quad f(x)=f(x')\Rightarrow x=x'.
\]
It is \emph{surjective} if every element of its specified codomain occurs as a value:
\[
\forall y\in F\ \exists x\in E,\quad f(x)=y.
\]
It is \emph{bijective} if both properties hold. A bijection pairs each input with
one output and each output with one input. Counting will use such pairings to
show that two finite sets have the same number of elements.

The function $x\mapsto x^2$ from $\R$ to $\R$ has neither property:
$1$ and $-1$ have the same image, and $-1$ has no preimage. Restricting its codomain
to $[0,\infty)$ makes it surjective. Restricting both domain and codomain to
$[0,\infty)$ makes it bijective. For injectivity there, $x^2=(x')^2$ gives
$(x-x')(x+x')=0$; either $x=x'$ or $x+x'=0$, and non-negativity makes the latter
case $x=x'=0$. Surjectivity is exactly the existence of non-negative square roots,
part of the later development from real completeness. The example distinguishes the
role of the domain from the role of the codomain.

Given $f:E\to F$ and $g:F\to G$, their \emph{composition} is
$g\circ f:E\to G$, defined by $(g\circ f)(x)=g(f(x))$.
Read the rightmost function first. For instance, if $f(x)=x+1$ and $g(x)=x^2$
on $\R$, then $g\circ f$ sends $x$ to $(x+1)^2$, whereas $f\circ g$ sends it
to $x^2+1$. The two disagree at $x=1$. Composition is associative because
$h(g(f(x)))$ is the value of both $(h\circ g)\circ f$ and $h\circ(g\circ f)$.
The \emph{identity function} $\operatorname{id}_E$ sends each $x\in E$ to itself.

An \emph{inverse function} for $f:E\to F$ is a function $g:F\to E$ satisfying
both $g\circ f=\operatorname{id}_E$ and $f\circ g=\operatorname{id}_F$.
The first equation undoes $f$ on inputs; the second undoes it on outputs.

\begin{theorem}[the inverse criterion]\label{thm:lang-inverse}
A function has an inverse function if and only if it is bijective.
The inverse, when it exists, is unique and is denoted $f^{-1}$.
\end{theorem}
\begin{proof}
Suppose $g$ is an inverse. If $f(x)=f(x')$, apply $g$ and use $g(f(x))=x$ to get
$x=x'$, so $f$ is injective. For each $y\in F$, the equation $f(g(y))=y$
provides a preimage, so $f$ is surjective.

Conversely, suppose $f$ is bijective. For each $y\in F$, surjectivity supplies an
$x\in E$ with $f(x)=y$, and injectivity makes that $x$ unique. Define $g(y)$ to
be this unique $x$. This gives $f(g(y))=y$; uniqueness also gives $g(f(x))=x$.
Thus $g$ is an inverse. If $h$ is another inverse, then for every $y\in F$,
$h(y)=h(f(g(y)))=g(y)$, hence $h=g$.
\end{proof}

For $f:\R\to\R$, $f(x)=3x-2$, solve $y=3x-2$ for $x$ to find the candidate
$g(y)=(y+2)/3$. Substituting into both compositions verifies the inverse.
Solving symbolically is a way to discover the formula, not a substitute for checking
its domains and the two identities. For a bijection, the inverse image of $B$
is also the image of $B$ under the inverse function, so the two uses of $f^{-1}(B)$ agree.

\begin{think}
The maps $f:\N\to\N$, $f(n)=n+1$, and $g:\N\to\N$, with $g(0)=0$ and
$g(n)=n-1$ for $n\ge1$, satisfy $g\circ f=\operatorname{id}_{\N}$.
Does $f\circ g=\operatorname{id}_{\N}$? What does this say about checking only
one of the two inverse identities?
\end{think}

\section{Relations, equivalence classes and partitions}

A \emph{relation} on a set $E$ is a subset $R\subseteq E\times E$.
Writing $x\mathrel R y$ means $(x,y)\in R$. Equality, $\le$ on $\R$, and
``has the same square as'' on $\R$ are relations with very different behaviour.
The useful properties have names:
\begin{itemize}
\item \emph{reflexive}: $x\mathrel R x$ for every $x\in E$;
\item \emph{symmetric}: $x\mathrel R y$ implies $y\mathrel R x$;
\item \emph{transitive}: $x\mathrel R y$ and $y\mathrel R z$ imply $x\mathrel R z$.
\end{itemize}
A relation satisfying all three is an \emph{equivalence relation}, usually written $\sim$.
It records a precise sense in which objects are to be treated as alike.
The relation $x\sim y$ defined by $x^2=y^2$ is an equivalence relation:
each property follows from the corresponding property of equality.
By contrast, $x\sim y$ defined by $|x-y|<1$ on $\R$ is reflexive and symmetric
but not transitive: $0$ is related to $3/4$ and $3/4$ to $3/2$, while $0$ is
not related to $3/2$. Here $|t|$ denotes $t$ for $t\ge0$ and $-t$ for $t<0$.

The \emph{equivalence class} of $x$ is
$[x]=\{y\in E:y\sim x\}$. For the square relation, $[2]=\{-2,2\}$
and $[0]=\{0\}$. The classes need not have equal sizes. A \emph{partition}
of $E$ is a collection of nonempty subsets, called its \emph{blocks}, whose union
is $E$ and whose distinct blocks are disjoint. The empty collection partitions
the empty set. A partition sorts objects into blocks without omissions or ambiguity.

\begin{theorem}[equivalence relations are partitions]\label{thm:lang-partition}
The distinct equivalence classes of an equivalence relation on $E$ form a partition
of $E$. Conversely, any partition defines an equivalence relation by declaring
two elements equivalent exactly when they lie in the same block; its classes are
the original blocks.
\end{theorem}
\begin{proof}
Reflexivity gives $x\in[x]$, so the classes are nonempty and cover $E$.
If $z\in[x]\cap[y]$, then $z\sim x$ and $z\sim y$. Symmetry and transitivity
give $x\sim y$. For any $t\in[x]$, transitivity now gives $t\sim y$,
so $[x]\subseteq[y]$; interchanging $x,y$ gives equality. Thus classes that meet
are equal, and distinct classes are disjoint.

Conversely, let a partition be given. Every element lies in exactly one block,
so it is related to itself. Sharing a block is symmetric. If $x,y$ share one
block and $y,z$ share another, those blocks meet at $y$, hence are equal;
then $x,z$ share a block. This proves transitivity. Finally, the elements
related to $x$ are exactly those in its block, proving the assertion about classes.
\end{proof}

This correspondence will return in several forms: grouping integers by remainders,
grouping outcomes by the information a random variable reveals, and grouping vectors
which a linear map sends to the same output. The definition supplies a common language
before any of those theories is developed.

\section{Order: comparable does not mean equal}

A \emph{partial order} on $E$ is a reflexive, transitive relation $\preceq$
which is \emph{antisymmetric}: if $x\preceq y$ and $y\preceq x$, then $x=y$.
Antisymmetry differs from symmetry. In an order, two-way comparison forces equality;
in an equivalence relation, it is expected even for distinct equivalent objects.
The usual $\le$ on $\R$ is an order. So is inclusion on $\mathcal P(E)$:
every set contains itself, inclusions compose, and two opposite inclusions give equality.
An order is \emph{total} if any two elements are comparable: $x\preceq y$ or
$y\preceq x$. Inclusion is usually only partial: $\{1\}$ and $\{2\}$ are incomparable.
The associated strict order $x\prec y$ means $x\preceq y$ and $x\ne y$.

Let $A$ be a subset of an ordered set. A \emph{least element} $a$ belongs to $A$
and satisfies $a\preceq x$ for every $x\in A$. A \emph{minimal element} belongs
to $A$ and has no strictly smaller element in $A$. Least implies minimal:
if $x\prec a$ with $x\in A$, leastness gives $a\preceq x$, contradicting
antisymmetry. The converse can fail in a partial order. Under inclusion,
$\{\{1\},\{2\}\}$ has two minimal elements and no least element.
In a total order the converse holds: if a minimal $a$ were not below some $x\in A$,
totality would give $x\prec a$, a contradiction.

A least element, if it exists, is unique: two such elements are below each other.
Replacing ``below'' by ``above'' defines \emph{greatest} and \emph{maximal} elements.
An \emph{upper bound} for $A$ is an element $u$ of the ambient ordered set satisfying
$x\preceq u$ for every $x\in A$; it need not belong to $A$. A \emph{least upper
bound}, or \emph{supremum}, is the least of its upper bounds. Reversing the order
defines a \emph{lower bound} and a \emph{greatest lower bound}, or \emph{infimum}.
For $(0,1)\subset\R$, $1$ is a supremum but not a greatest element.
Indeed, $1$ bounds the interval above; if $u<1$, the number
$(\max\{u,0\}+1)/2$ lies in $(0,1)$ and exceeds $u$, so $u$ is not an upper bound.
Here $\max\{u,0\}$ is the larger of the two real numbers $u,0$.
Completeness asserts that every
nonempty real set bounded above has a supremum. The sequences chapter will show
how this assumption produces limits and existence theorems.

\section{Choosing a proof method}

The form of a statement often suggests the first line of a proof. To prove
``for every $x\in E$, $P(x)$'', let $x$ be an arbitrary element of $E$ and
establish $P(x)$ without adding special restrictions. To prove ``there exists'',
construct a witness or prove that a witness must exist. To prove an equivalence,
write down both directions. To prove equality of sets, compare membership in
both directions. These are obligations imposed by the words, before any clever idea enters.

\subsection*{Direct proof and proof by cases}

A \emph{direct proof} starts from the hypotheses and reaches the conclusion.
The even-square argument was direct. Here is another: if $a<b$ are real numbers,
there is a real number strictly between them. Take $m=(a+b)/2$. Then
$m-a=(b-a)/2>0$ and $b-m=(b-a)/2>0$, so $a<m<b$.
The proof supplies a witness and verifies exactly the two required inequalities.

When the hypotheses give alternatives, a proof can handle them separately. For
real $x$, $|x|\ge x$: if $x\ge0$, equality holds; if $x<0$, then
$|x|=-x>x$. A proof by cases must cover every allowed case. Dividing into
$x>0$ and $x<0$ would omit $x=0$.

\subsection*{Contrapositive}

The implication $P\Rightarrow Q$ is equivalent to its \emph{contrapositive}
$\neg Q\Rightarrow\neg P$. Both fail in exactly the case $P$ true and $Q$ false.
This offers a different starting point, without changing the theorem.
For example, if $x$ is irrational, then $x+1$ is irrational. Prove the
contrapositive: if $x+1=a/b$ for integers $a,b$ with $b\ne0$, then
$x=(a-b)/b$ is rational. The assumption we used directly is the negation of
the desired conclusion, and the conclusion reached is the negation of the
original assumption. This is not the converse ``if $x+1$ is irrational, then
$x$ is irrational'', even though that converse happens also to be true.

\subsection*{Contradiction}

A \emph{proof by contradiction} assumes the hypotheses and the negation of
the desired conclusion, then derives an impossibility. Classical logic then
rules out that negation. There is no smallest positive real number: if $a>0$
were smallest, then $a/2>0$ and $a/2<a$, contradicting its supposed minimality.
The contradiction is a pair of incompatible statements, not merely an unexpected answer.

The two indirect methods overlap, but their targets differ. A contrapositive proof
of $P\Rightarrow Q$ proves a new implication $\neg Q\Rightarrow\neg P$.
A contradiction proof may begin with $P$ and $\neg Q$ together and derive any
impossibility. Use whichever makes the available information productive; do not
turn a direct calculation into a contradiction proof merely by adding ``suppose not''.

\subsection*{Counterexamples, existence and uniqueness}

To disprove $\forall x\in E,\ P(x)$, give one $x\in E$ and show that $P(x)$ fails.
For an implication, the example must satisfy its hypothesis. To disprove
``every positive real $x$ has $x^2>x$'', $x=-1$ is irrelevant, but $x=1/2$ works:
$1/4<1/2$. Counterexamples often live at $0$, at an endpoint, at equal parameters,
or in a small structure where an informal picture breaks down.

An \emph{existence and uniqueness} claim has two parts. The notation
$\exists!x\in E,\ P(x)$ abbreviates
\[
\exists x\in E:\quad P(x)\ \hbox{and}\
\forall y\in E,\ (P(y)\Rightarrow y=x).
\]
Finding a candidate proves neither part by itself. For the equation $3x+2=8$
over $\R$, rearrangement suggests $x=2$. Substitution proves existence.
If $y$ is any solution, $3y+2=8$ gives $y=2$, proving uniqueness.
In harder problems an existence proof may identify no explicit formula; the
intermediate value theorem will provide an important example. Whenever the book
claims uniqueness, look for the step that excludes a second solution.

\begin{reminder}[errors that survive a convincing calculation]
An argument can fail by proving the converse, choosing an existential witness
before the data it must depend on, checking only examples, assuming the conclusion,
or using an expression outside its domain. Dividing by a quantity requires knowing
it is nonzero; squaring an equation may introduce extraneous solutions; taking a
square root produces the non-negative root, so $\sqrt{x^2}=|x|$.
Finally, write each ``therefore'' as a small question: does the information on
the page really force the next assertion?
\end{reminder}

\section{Proof workshop: when does an image preserve an intersection?}

The elementary-looking question
\[
f(A\cap B)\stackrel{?}=f(A)\cap f(B)
\]
combines definitions, quantifiers, counterexamples and a useful characterisation.
Here $f:E\to F$ is a function and $A,B\subseteq E$. An element of the left side
comes from one input which belongs to both sets. An element of the right side
comes from an input in $A$ and an input in $B$; the two inputs need not agree.
That is where to experiment.

\subsection*{A worked experiment and a conjecture}

For $f:\R\to\R$, $f(x)=x^2$, choose $A=\{-1\}$ and $B=\{1\}$.
Then $A\cap B=\varnothing$, so $f(A\cap B)=\varnothing$, while
$f(A)\cap f(B)=\{1\}$. Equality fails. Yet the left side seems always
contained in the right: any common input supplies a common output.
The failure used two different inputs with the same output. This suggests that
injectivity may be exactly the missing hypothesis.

\begin{think}
Prove the inclusion that always holds. Then suppose $f$ is injective and
try to reverse it. Which two witnesses does membership in $f(A)\cap f(B)$ supply,
and what allows you to identify them?
\end{think}

\subsection*{Mission and staged hints}

Prove both assertions:
\begin{enumerate}[label=(\alph*)]
\item For every function and every $A,B\subseteq E$,
$f(A\cap B)\subseteq f(A)\cap f(B)$; equality holds if $f$ is injective.
\item If equality holds for \emph{every} pair of subsets $A,B\subseteq E$,
then $f$ is injective.
\end{enumerate}
\emph{First hint.} Start with an arbitrary element of the set whose inclusion
you want to prove, and translate ``image'' into a witness.
\emph{Second hint.} For (b), you are allowed to choose the sets. If $f(x)=f(x')$,
what do $A=\{x\}$ and $B=\{x'\}$ reveal?

\subsection*{How to find it}

The first direction follows the definitions forwards. The failed reverse direction
reveals a specific obstruction: two witnesses can disagree. Injectivity removes
that obstruction. To show the hypothesis is necessary, make the obstruction as
small as possible by using singletons. This is a common discovery pattern:
attempt a proof, locate the exact failing step, then either repair its hypothesis
or turn the failure into a counterexample.

\subsection*{Write-up}

Let $y\in f(A\cap B)$. There is $x\in A\cap B$ with $f(x)=y$.
This $x$ belongs to each of $A,B$, so $y\in f(A)\cap f(B)$.
Now assume $f$ injective and let $y\in f(A)\cap f(B)$. There are $a\in A$
and $b\in B$ with $f(a)=y=f(b)$. Injectivity gives $a=b$, hence
$a\in A\cap B$ and $y\in f(A\cap B)$. This proves equality.

Conversely, assume equality for every $A,B\subseteq E$. If $f(x)=f(x')=y$,
take $A=\{x\}$ and $B=\{x'\}$. Then $y\in f(A)\cap f(B)$, so the
assumed equality gives an element of $A\cap B$. Two singleton sets meet only
when their elements agree. Thus $x=x'$, proving injectivity. The argument also
covers $E=\varnothing$: the injectivity condition and all the inclusions are vacuous.

\subsection*{A tempting invalid proof}

``If $y\in f(A)\cap f(B)$, write $y=f(x)$ with $x\in A$ and $x\in B$.''
The single letter $x$ has silently identified two independent existential choices.
The correct translation initially supplies $a\in A$ and $b\in B$.
Give different witnesses different names until an argument proves they coincide.
This small writing discipline prevents a large class of errors.

\section{Exercises}

\begin{exercise}[$\star$]\label{ex:lang-implications}
For real $x$, decide which implications hold for every $x$ and give a counterexample
to each that fails:
(a) $x=2\Rightarrow x^2=4$;
(b) $x^2=4\Rightarrow x=2$;
(c) $x>2\Rightarrow x^2>4$;
(d) $x^2>4\Rightarrow x>2$.
Give a necessary and sufficient condition for $x^2>4$.
\end{exercise}

\begin{exercise}[$\star$]\label{ex:lang-negation}
Negate each statement, preserving its domain:
(a) every $x\in\R$ satisfies $x^2\ge x$;
(b) some $x\in\R$ satisfies $0<x<1$ and $x^2=x$;
(c) for every $x\in\R$ there is $y\in\R$ with $y^2=x$.
Decide whether the original statements are true.
\end{exercise}

\begin{exercise}[$\star$]\label{ex:lang-quantifiers}
Decide which statements are true and justify the answers:
\[
\forall x\in\R\ \exists y\in\R:\ y>x,
\qquad
\exists y\in\R\ \forall x\in\R:\ y>x.
\]
Write their negations and explain which choices may depend on which variables.
\end{exercise}

\begin{exercise}[$\star$]\label{ex:lang-setalgebra}
Prove $A\setminus(B\cup C)=(A\setminus B)\cap(A\setminus C)$ by comparing
membership. Is $A\setminus(B\cap C)=(A\setminus B)\cap(A\setminus C)$ always true?
If it fails, give a counterexample and a corrected equality.
\end{exercise}

\begin{exercise}[$\star$]\label{ex:lang-powerset}
List $\mathcal P(\{\varnothing,\{1\}\})$. Which of the following hold:
$\varnothing\in\mathcal P(A)$, $\varnothing\subseteq\mathcal P(A)$,
$A\in\mathcal P(A)$, and $A\subseteq\mathcal P(A)$, for an arbitrary set $A$?
Give a counterexample to any assertion which need not hold.
\end{exercise}

\begin{exercise}[$\star\star$]\label{ex:lang-domain}
For $E=\R\setminus\{0\}$, consider $f:E\to\R$, $f(x)=1/x$.
Prove that it is injective but not surjective. Change only its codomain to make
it bijective, and find and verify the inverse. Explain why the codomain matters.
\end{exercise}

\begin{exercise}[$\star\star$]\label{ex:lang-inverseimage}
For a function $f:E\to F$ and $A\subseteq E$, $B\subseteq F$, prove
\[
A\subseteq f^{-1}(f(A)),\qquad f(f^{-1}(B))=B\cap f(E).
\]
Give an example where the first inclusion is strict. Characterise the functions
for which it is equality for every $A\subseteq E$.
\end{exercise}

\begin{exercise}[$\star\star$]\label{ex:lang-complement}
For $f:E\to F$, prove $f(E\setminus A)\supseteq f(E)\setminus f(A)$.
Give a counterexample to equality. Prove equality when $f$ is injective.
Explain why replacing $f(E)$ by $F$ would also require surjectivity.
\end{exercise}

\begin{exercise}[$\star\star$]\label{ex:lang-composition}
Let $f:E\to F$ and $g:F\to G$. Prove:
(a) if both are injective, then $g\circ f$ is injective;
(b) if both are surjective, then $g\circ f$ is surjective;
(c) if $g\circ f$ is injective, then $f$ is injective;
(d) if $g\circ f$ is surjective, then $g$ is surjective.
Give examples showing that the other factor need not have the same property in (c) or (d).
\end{exercise}

\begin{exercise}[$\star\star$]\label{ex:lang-selfinverse}
Let $E=\R\setminus\{1\}$. Show that $f(x)=(x+1)/(x-1)$ defines a function
$E\to E$ and satisfies $f\circ f=\operatorname{id}_E$.
Deduce bijectivity. Check every denominator before cancelling it.
\end{exercise}

\begin{exercise}[$\star\star$]\label{ex:lang-equivalence}
Define $x\sim y$ on $\R$ by $x-y\in\Z$. Prove this is an equivalence relation.
Describe $[0]$ and $[1/2]$, and decide whether $[1/2]=[3/2]$.
Can two distinct equivalence classes share an element?
\end{exercise}

\begin{exercise}[$\star\star$]\label{ex:lang-partitions}
Suppose $f:E\to F$. Define $x\sim y$ if $f(x)=f(y)$.
Prove that $\sim$ is an equivalence relation and that its classes are exactly
the nonempty fibres $f^{-1}(\{y\})$, for $y\in f(E)$.
What does injectivity say about these classes?
\end{exercise}

\begin{exercise}[$\star\star$]\label{ex:lang-order}
Order $\mathcal A=\{\{1\},\{2\},\{1,2\}\}$ by inclusion. Find all its minimal,
maximal, least and greatest elements. Prove that any least element of any partially
ordered set is unique. Is the empty set an upper or a lower bound for $\mathcal A$
in the larger ordered set $\mathcal P(\{1,2\})$?
\end{exercise}

\begin{exercise}[$\star\star$]\label{ex:lang-rational}
Let $r\in\Q$ and $x\in\R$. Prove that $x+r$ is rational if and only if $x$ is rational.
If also $r\ne0$, prove that $rx$ is rational if and only if $x$ is rational.
State the corresponding consequences for irrational numbers, and explain the role of $r\ne0$.
\end{exercise}

\begin{exercise}[$\star\star$]\label{ex:lang-existence}
For each $a\in(0,1)$, prove that there is a unique $x>0$ satisfying
$x/(1+x)=a$. Explain separately how the assumptions give an admissible candidate
and how the equation gives uniqueness.
\end{exercise}

\begin{exercise}[$\star$]\label{ex:lang-falseproof}
Find the first invalid step in this purported proof that $1=2$. Let $a=b=1$.
Then $a^2=ab$, hence $a^2-b^2=ab-b^2$, hence
$(a-b)(a+b)=b(a-b)$, hence $a+b=b$, hence $2=1$.
Why do the valid lines before the error fail to make the conclusion plausible evidence?
\end{exercise}

\begin{exercise}[$\star\star$]\label{ex:lang-midpoint}
For real $a<b$, prove that there is exactly one real $x$ with
$|x-a|=|x-b|$, and determine it. Give one proof by cases according to the position
of $x$, and another by squaring. Explain why squaring is reversible in this argument.
\end{exercise}

\begin{exercise}[$\star\star\star$]\label{ex:lang-retraction}
Let $f:E\to E$ satisfy $f\circ f=f$.
(a) Prove $f(E)=\{y\in E:f(y)=y\}$.
(b) Prove that each class of the relation $x\sim x'$ defined by $f(x)=f(x')$
contains exactly one fixed point of $f$.
(c) If such an $f$ is injective or surjective, prove $f=\operatorname{id}_E$.
Give an example with $f\ne\operatorname{id}_E$ when neither property is required.
\end{exercise}

\section{Hints}

\begin{hint}{ex:lang-implications}
Test the negative root in (b), and a number below $-2$ in (d). Factor $x^2-4$.
\end{hint}
\begin{hint}{ex:lang-negation}
Change each quantifier, then negate its predicate. The negation of an ``and'' is an ``or''.
\end{hint}
\begin{hint}{ex:lang-quantifiers}
For a given $x$, try $x+1$. For a proposed universal $y$, choose $x$ using that $y$.
\end{hint}
\begin{hint}{ex:lang-setalgebra}
Use the phrases ``outside both'' and ``outside at least one''. Test sets with one element.
\end{hint}
\begin{hint}{ex:lang-powerset}
Your base set has two elements even though one is itself empty. Membership in
$\mathcal P(A)$ means being a subset of $A$.
\end{hint}
\begin{hint}{ex:lang-domain}
Which real value can a reciprocal never take? To solve $1/x=y$, you must know $y\ne0$.
\end{hint}
\begin{hint}{ex:lang-inverseimage}
Translate each membership into a witness. For the characterisation, use singleton $A$.
\end{hint}
\begin{hint}{ex:lang-complement}
A value in $f(E)$ but not in $f(A)$ has a preimage, and none of its preimages is in $A$.
For strictness, let two different inputs share a value.
\end{hint}
\begin{hint}{ex:lang-composition}
Write out each definition. For failed converses, a singleton intermediate image can hide
bad behaviour elsewhere in $F$.
\end{hint}
\begin{hint}{ex:lang-selfinverse}
First solve the impossible equation $(x+1)/(x-1)=1$. Then compute $f(x)-1$.
\end{hint}
\begin{hint}{ex:lang-equivalence}
Use $0\in\Z$, closure under negation and closure under addition. Classes consist of
integer translates.
\end{hint}
\begin{hint}{ex:lang-partitions}
Equality of outputs inherits reflexivity, symmetry and transitivity. Fix the output of one input.
\end{hint}
\begin{hint}{ex:lang-order}
Draw the three sets with arrows for strict inclusion. A bound in the larger ordered set
need not belong to $\mathcal A$.
\end{hint}
\begin{hint}{ex:lang-rational}
Use closure of rational numbers under addition and subtraction, and under division by
a nonzero rational. Write fractions to justify closure if needed.
\end{hint}
\begin{hint}{ex:lang-existence}
Solve the equation backwards, then check the candidate forwards. The sign of $1-a$ matters.
\end{hint}
\begin{hint}{ex:lang-falseproof}
Which factor was cancelled, and what is its value under the original assumptions?
\end{hint}
\begin{hint}{ex:lang-midpoint}
Outside $[a,b]$, one distance is strictly larger. Inside, the equation is linear.
For the second proof, use the difference of two squares.
\end{hint}
\begin{hint}{ex:lang-retraction}
Apply the defining equation to an element of the image. In the class of $x$, try the representative $f(x)$.
\end{hint}

\section{Solutions}

\begin{solution}{ex:lang-implications}
\idea Distinguish a permitted input from an arbitrary input; squares lose the sign of a number.
\writeup (a) is true by substitution. (b) fails at $x=-2$. For (c), $x>2$ gives
$x-2>0$ and $x+2>0$, so $x^2-4=(x-2)(x+2)>0$. (d) fails at $x=-3$.
The necessary and sufficient condition is $x<-2$ or $x>2$. If $-2\le x\le2$,
then $|x|\le2$ and $x^2\le4$; outside that interval, $|x|>2$ gives $x^2>4$.
\end{solution}

\begin{solution}{ex:lang-negation}
\idea Keep the domain fixed and negate from the outside in; then test the simplest candidates.
\writeup (a) negates to $\exists x\in\R,\ x^2<x$; the original is false at $x=1/2$.
(b) negates to
$\forall x\in\R,\ x\le0\ \hbox{or}\ x\ge1\ \hbox{or}\ x^2\ne x$.
The original is false: $x^2=x$ factors as $x(x-1)=0$, giving only $0,1$, both excluded.
(c) negates to $\exists x\in\R\ \forall y\in\R,\ y^2\ne x$.
The original is false, since $x=-1$ cannot be a real square.
\end{solution}

\begin{solution}{ex:lang-quantifiers}
\idea Ask what information is available at the moment $y$ is chosen.
\writeup The first is true: given $x$, choose $y=x+1>x$. Its negation is
$\exists x\in\R\ \forall y\in\R,\ y\le x$.
The second is false: if $y$ were such a witness, taking $x=y$ would require $y>y$.
Its negation, $\forall y\in\R\ \exists x\in\R,\ y\le x$, is true with $x=y$.
In the first original statement $y$ may depend on $x$; in the second it must
work for all subsequent choices of $x$.
\end{solution}

\begin{solution}{ex:lang-setalgebra}
\idea Difference translates into membership together with non-membership; use De Morgan's rules.
\writeup An $x$ lies in $A\setminus(B\cup C)$ exactly when it belongs to $A$
and belongs to neither $B$ nor $C$. This is equivalent to belonging to both
$A\setminus B$ and $A\setminus C$, giving the first equality.
The proposed second equality fails for $A=B=\{1\}$ and $C=\varnothing$:
its left side is $\{1\}$ and its right side is empty. The corrected equality is
$A\setminus(B\cap C)=(A\setminus B)\cup(A\setminus C)$.
Indeed, an $x\in A$ fails to belong to $B\cap C$ exactly when at least one
of its memberships in $B,C$ fails.
\end{solution}

\begin{solution}{ex:lang-powerset}
\idea Treat the two elements as named objects first; only then inspect what they contain.
\writeup The power set is
\[
\{\varnothing,\{\varnothing\},\{\{1\}\},\{\varnothing,\{1\}\}\}.
\]
The first three assertions always hold: the empty set is a subset of every $A$,
it is a subset of every $\mathcal P(A)$, and $A$ is a subset of itself, so
$A\in\mathcal P(A)$. The last need not hold. Take $A=\{\{\varnothing\}\}$:
its element $\{\varnothing\}$ is not a
subset of $A$, since $\varnothing\notin A$.
\end{solution}

\begin{solution}{ex:lang-domain}
\idea Solve equality of outputs for injectivity, then list the missing outputs.
\writeup If $1/x=1/x'$ with nonzero $x,x'$, multiplication by $xx'$ gives
$x=x'$. Thus $f$ is injective. The value $0$ never occurs, so $f:E\to\R$
is not surjective. With codomain $E$, it is bijective: for every $y\ne0$,
$x=1/y\in E$ gives $f(x)=y$. The inverse is itself, because
$1/(1/x)=x$ for every nonzero $x$. Specifying $E$ as codomain removes exactly
the unattainable output, without changing the values of the rule.
\end{solution}

\begin{solution}{ex:lang-inverseimage}
\idea A forward trip followed by an inverse image can collect additional inputs sharing an output.
\writeup If $x\in A$, then $f(x)\in f(A)$, hence $x\in f^{-1}(f(A))$.
An element $y$ belongs to $f(f^{-1}(B))$ precisely when some $x\in E$ has
$f(x)=y\in B$. This says exactly $y\in B\cap f(E)$.
For $f:\R\to\R$, $f(x)=x^2$, $A=\{1\}$, the first inverse image is
$\{-1,1\}$, so the inclusion is strict.
If $f$ is injective and $x\in f^{-1}(f(A))$, choose $a\in A$ with
$f(x)=f(a)$; then $x=a\in A$, proving equality. Conversely, if equality
holds for every $A$ and $f(x)=f(x')$, take $A=\{x\}$: then
$x'\in f^{-1}(f(A))=A$, hence $x'=x$. Thus the exact condition is injectivity.
\end{solution}

\begin{solution}{ex:lang-complement}
\idea An output absent from $f(A)$ cannot have any preimage in $A$; an output present
in $f(A)$ can still have another preimage outside $A$.
\writeup If $y\in f(E)\setminus f(A)$, choose $x\in E$ with $f(x)=y$.
Since $y\notin f(A)$, $x\notin A$, so $y\in f(E\setminus A)$.
For the square map and $A=\{1\}$, $1$ lies in $f(E\setminus A)$ via $x=-1$,
but not in $f(E)\setminus f(A)$, so equality can fail.
If $f$ is injective and $y=f(x)$ with $x\notin A$, there cannot be an
$a\in A$ with $f(a)=y$, as that would force $a=x$. Thus $y\notin f(A)$,
giving the reverse inclusion. Finally, if equality with $F\setminus f(A)$
is required for every $A$, take $A=\varnothing$: it says $f(E)=F$.
Under injectivity and surjectivity together, the equality follows from the one proved.
\end{solution}

\begin{solution}{ex:lang-composition}
\idea Trace equal outputs backwards for injectivity and a desired output backwards for surjectivity.
\writeup (a) If $g(f(x))=g(f(x'))$, injectivity of $g$ gives $f(x)=f(x')$,
then injectivity of $f$ gives $x=x'$. (b) Given $z\in G$, first choose
$y\in F$ with $g(y)=z$, then $x\in E$ with $f(x)=y$; thus $g(f(x))=z$.
(c) If $f(x)=f(x')$, applying $g$ and using injectivity of $g\circ f$ gives $x=x'$.
(d) Given $z\in G$, surjectivity of $g\circ f$ supplies $x$ with
$z=g(f(x))$, hence a preimage under $g$.
For both failed converses use $E=G=\{0\}$, $F=\{0,1\}$, $f(0)=0$ and
$g(0)=g(1)=0$. The composite is a bijection of a singleton, but $g$ is
not injective and $f$ is not surjective.
\end{solution}

\begin{solution}{ex:lang-selfinverse}
\idea Establish that the second application is legal before simplifying it.
\writeup For $x\ne1$, $f(x)$ is defined. If $f(x)=1$, then $x+1=x-1$,
which would give $1=-1$. Thus $f(x)\in E$. Moreover
\[
f(x)+1=\frac{2x}{x-1},\qquad f(x)-1=\frac2{x-1}\ne0,
\]
and consequently $f(f(x))=(2x/(x-1))/(2/(x-1))=x$.
This single identity supplies both inverse identities because the two maps
are the same. The inverse criterion gives bijectivity.
\end{solution}

\begin{solution}{ex:lang-equivalence}
\idea Differences add along a chain of equivalent elements.
\writeup Since $x-x=0\in\Z$, the relation is reflexive. If $x-y\in\Z$,
then $y-x=-(x-y)\in\Z$, proving symmetry. If also $y-z\in\Z$, then
$x-z=(x-y)+(y-z)\in\Z$, proving transitivity. The classes are
$[0]=\Z$ and $[1/2]=\{1/2+k:k\in\Z\}$.
As $3/2-1/2=1\in\Z$, the latter is also $[3/2]$.
Two classes which share an element are equal by Theorem~\ref{thm:lang-partition},
so distinct classes cannot meet.
\end{solution}

\begin{solution}{ex:lang-partitions}
\idea A fibre records exactly the information preserved by a function.
\writeup Reflexivity, symmetry and transitivity follow respectively from
$f(x)=f(x)$, reversal of equality, and transitivity of equality.
For $x\in E$, its class consists of all $x'\in E$ satisfying $f(x')=f(x)$,
so it is the nonempty fibre $f^{-1}(\{f(x)\})$.
Conversely every $y\in f(E)$ equals $f(x)$ for some $x$, so its fibre is that
class. Injectivity holds exactly when all these classes have one element:
each contains its defining input, and injectivity forbids a second one.
For $E=\varnothing$ there are no classes, and this condition is vacuously true.
\end{solution}

\begin{solution}{ex:lang-order}
\idea A least element must be below every other element; minimality only rules out smaller ones.
\writeup The minimal elements are $\{1\}$ and $\{2\}$. Neither contains a
strictly smaller member of $\mathcal A$. Neither is below the other, so
there is no least element. The only maximal element is $\{1,2\}$, which
is also greatest because it contains every member of $\mathcal A$.
If $a,b$ are both least in any partial order, then $a\preceq b$ and
$b\preceq a$, hence $a=b$ by antisymmetry. In the larger power-set order,
$\varnothing$ is a lower bound for $\mathcal A$ because it is a subset
of all its members. It is not an upper bound, since it does not contain $\{1\}$.
\end{solution}

\begin{solution}{ex:lang-rational}
\idea The reverse implications subtract or divide by the fixed rational parameter.
\writeup Rational numbers are closed under addition, subtraction and multiplication,
and division by a nonzero rational: these follow by putting fractions over common
nonzero denominators and, for division, multiplying by the reciprocal of a fraction
with nonzero numerator. Thus $x\in\Q$ gives $x+r\in\Q$, while
$x+r\in\Q$ gives $x=(x+r)-r\in\Q$.
For $r\ne0$, $x\in\Q$ gives $rx\in\Q$, and $rx\in\Q$ gives
$x=(rx)/r\in\Q$. Negating the equivalences says that adding a rational,
or multiplying by a nonzero rational, preserves irrationality in both directions.
The restriction is essential: $0x=0$ is rational for every real $x$.
\end{solution}

\begin{solution}{ex:lang-existence}
\idea Algebra suggests the candidate; the signs of the parameters verify it lies in the required domain.
\writeup Given $0<a<1$, set $x=a/(1-a)$. Its numerator and denominator
are positive, hence $x>0$, and $1+x=1/(1-a)>0$. Therefore
$x/(1+x)=a$. This proves existence.
If $y>0$ also satisfies $y/(1+y)=a$, the nonzero denominator permits
multiplication to obtain $y=a(1+y)$, hence $(1-a)y=a$. Since $1-a>0$,
division gives $y=a/(1-a)=x$. This proves uniqueness.
\end{solution}

\begin{solution}{ex:lang-falseproof}
\idea Cancellation is division, even when no division sign is written.
\writeup The step from $(a-b)(a+b)=b(a-b)$ to $a+b=b$ cancels $a-b$.
But $a=b=1$, so this factor is zero. The preceding equation is simply
$0=0$, which places no restriction on $a+b$ versus $b$. Correct lines
which discard information cannot validate a later attempt to recover that information.
All subsequent numerical conclusions rely on the illegal cancellation.
\end{solution}

\begin{solution}{ex:lang-midpoint}
\idea A point outside the interval cannot be equidistant from its two distinct endpoints.
\writeup If $x<a$, then $|x-b|=b-x>a-x=|x-a|$; if $x>b$, then
$|x-a|=x-a>x-b=|x-b|$. For $a\le x\le b$, equality becomes
$x-a=b-x$, whose unique solution is $x=(a+b)/2$. This point lies strictly
between $a$ and $b$ and satisfies the equation, proving existence and uniqueness.
Alternatively, both absolute values are non-negative, so they are equal
if and only if their squares are equal. The equation is therefore equivalent to
\[
0=(x-a)^2-(x-b)^2=(b-a)(2x-a-b).
\]
Since $b-a>0$, it is equivalent to $x=(a+b)/2$.
Squaring is reversible here because two non-negative numbers with equal squares
are equal: their squared difference factors as their difference times their sum;
if the sum is zero both are zero, and otherwise the difference must vanish.
\end{solution}

\begin{solution}{ex:lang-retraction}
\idea The equation says that after one application, applying $f$ again changes nothing.
\writeup (a) If $y\in f(E)$, write $y=f(x)$; then
$f(y)=f(f(x))=f(x)=y$, so $y$ is fixed. Conversely, if $f(y)=y$, then
$y$ is the image of itself, hence belongs to $f(E)$.
(b) In the class of any $x\in E$, the element $f(x)$ is fixed by (a) and
has the same output as $x$, so it belongs to that class. If fixed points
$y,z$ lie in the same class, then $y=f(y)=f(z)=z$; thus this representative is unique.
(c) If $f$ is injective, equality $f(f(x))=f(x)$ gives $f(x)=x$ for every $x$.
If $f$ is surjective, (a) says every element of $E$ is fixed, giving the same conclusion.
For an example, take $E=\{0,1\}$ and $f(0)=f(1)=0$.
It satisfies $f\circ f=f$ but is neither injective nor surjective onto $E$,
and differs from the identity. For the empty set all conclusions remain valid.
\end{solution}

\section*{Challenges and onward connections}
\begin{challenge}
\textbf{1. Recover a partition from information.} Given a partition of a set $E$,
map each element to its block. Prove that this map is surjective onto the set of
blocks and that its equal-output relation recovers the original partition.
\par\smallskip
\textbf{2. An order on subsets.} Prove that intersection is the greatest lower
bound and union the least upper bound of two subsets under inclusion.
\par\smallskip
\textbf{3. A larger infinity.} Can a function $f:E\to\mathcal P(E)$ ever be
surjective? Investigate the subset $D=\{x\in E:x\notin f(x)\}$.
\end{challenge}

\noindent\textbf{Notes.} In the first challenge, the unique block containing an
element makes the map well defined. Every block is nonempty, so every block is
hit, and equality of blocks is exactly the original relation. In the second,
$A\cap B$ is a subset of both; every common subset is contained in it.
$A\cup B$ contains both; every common superset contains it.
For the third, if $f(a)=D$, then $a\in D$ is equivalent to $a\notin f(a)$,
which is equivalent to $a\notin D$, an impossibility. Thus no such function is
surjective, even when $E$ is empty. This diagonal argument is a proved glimpse
of the theory of infinite sets developed further in Volume~6; it is not a
claim that all infinite sets can be compared by finite counting.

\section*{Reading and the next step}
For further practice in the language and organisation of proofs, see Daniel J. Velleman,
\emph{How to Prove It: A Structured Approach}, 3rd edition, Cambridge University Press,
2019; and Richard Hammack, \emph{Book of Proof}, 3rd edition, 2018, freely available
from the author's website. These are supplementary proof texts, not prerequisites.
The next chapter adds induction and minimal counterexamples: methods for passing
from a precise claim about one natural number to a theorem about them all.
